% ============================================================
%  DOCUMENTO LaTeX — Apuntes Universitarios
%  Materia: Derechos Reales — El Derecho Real de Superficie
%  Autor: Isaac Edgar Camacho Ocampo
%  Institución: Facultad de Derecho — Universidad de Buenos Aires (UBA)
%  Generado automáticamente a partir de apuntes de clase
% ============================================================

\documentclass[12pt,oneside]{book}

% ── METADATOS ──────────────────────────────────────────────
\newcommand{\universidad}{Universidad de Buenos Aires}
\newcommand{\facultad}{Facultad de Derecho}
\newcommand{\carrera}{Ciclo Superior --- Cátedra de Derechos Reales}
\newcommand{\materia}{Derechos Reales}
\newcommand{\titulo}{El Derecho Real de Superficie}
\newcommand{\autor}{Isaac Edgar Camacho Ocampo}
\newcommand{\anio}{2024}

% ── PAQUETES ───────────────────────────────────────────────
\usepackage[utf8]{inputenc}
\usepackage[T1]{fontenc}
\usepackage[spanish,es-nodecimaldot]{babel}

\usepackage[
    top=2.5cm,
    bottom=2.5cm,
    left=3cm,
    right=2.5cm,
    headheight=15pt
]{geometry}

\usepackage{amsmath}
\usepackage{amssymb}
\usepackage{mathtools}

\usepackage{graphicx}
\usepackage{float}
\usepackage{caption}
\usepackage{subcaption}

\usepackage{booktabs}
\usepackage{array}
\usepackage{longtable}
\usepackage{tabularx}

\usepackage{enumitem}

\usepackage{xcolor}
\usepackage[most]{tcolorbox}

\usepackage{fancyhdr}

\usepackage{ragged2e}

\usepackage[
    colorlinks=true,
    linkcolor=blue!50!black,
    citecolor=green!40!black,
    urlcolor=blue!60!black,
    pdftitle={\titulo},
    pdfauthor={\autor},
    pdfsubject={Apuntes universitarios},
    bookmarks=true
]{hyperref}

% ── RUTAS DE IMÁGENES ──────────────────────────────────────
\graphicspath{
    {./img/}
    {./graficos/}
    {./figuras/}
}

% ── CONFIGURACIÓN ──────────────────────────────────────────
\setcounter{tocdepth}{2}

\setlength{\parindent}{0pt}
\setlength{\parskip}{0.5em}

\setlist[itemize]{topsep=0.3em, itemsep=0.2em}
\setlist[enumerate]{topsep=0.3em, itemsep=0.2em}

\clubpenalty=10000
\widowpenalty=10000

% ── ENCABEZADOS ────────────────────────────────────────────
\pagestyle{fancy}
\fancyhf{}
\fancyhead[L]{\nouppercase{\leftmark}}
\fancyhead[R]{\materia}
\fancyfoot[C]{\thepage}
\renewcommand{\headrulewidth}{0.4pt}
\renewcommand{\footrulewidth}{0pt}

\fancypagestyle{plain}{
    \fancyhf{}
    \fancyfoot[C]{\thepage}
    \renewcommand{\headrulewidth}{0pt}
}

% ── COMANDOS MATEMÁTICOS ───────────────────────────────────
\newcommand{\abs}[1]{\left|#1\right|}
\newcommand{\norm}[1]{\left\lVert#1\right\rVert}
\newcommand{\vect}[1]{\mathbf{#1}}
\newcommand{\deriv}[2]{\frac{d#1}{d#2}}
\newcommand{\pderiv}[2]{\frac{\partial#1}{\partial#2}}

% ── CAJAS ACADÉMICAS ───────────────────────────────────────
\newtcolorbox{definition}[1]{
    enhanced, breakable,
    title={Definición: #1},
    fonttitle=\bfseries,
    colback=blue!3,
    colframe=blue!50!black,
    boxrule=0.8pt, arc=2mm
}

\newtcolorbox{important}{
    enhanced, breakable,
    title={Importante},
    fonttitle=\bfseries,
    colback=yellow!8,
    colframe=orange!70!black,
    boxrule=0.8pt, arc=2mm
}

\newtcolorbox{example}[1]{
    enhanced, breakable,
    title={Caso práctico: #1},
    fonttitle=\bfseries,
    colback=green!4,
    colframe=green!40!black,
    boxrule=0.8pt, arc=2mm
}

\newtcolorbox{formula}[1]{
    enhanced, breakable,
    title={Artículo / Regla: #1},
    fonttitle=\bfseries,
    colback=purple!4,
    colframe=purple!50!black,
    boxrule=0.8pt, arc=2mm
}

\newtcolorbox{note}{
    enhanced, breakable,
    title={Nota},
    fonttitle=\bfseries,
    colback=white,
    colframe=gray!50,
    boxrule=0.6pt, arc=2mm
}

% ============================================================
\begin{document}

% ── PORTADA ─────────────────────────────────────────────────
\begin{titlepage}
\thispagestyle{empty}
\begin{center}

\vspace*{1cm}

{\large \textsc{\universidad}}

\vspace{0.2cm}

{\large \textsc{\facultad}}

\vspace{0.2cm}

{\large \carrera}

\vspace{3cm}

{\Huge \textbf{\materia}}

\vspace{1cm}

{\LARGE \titulo}

\vspace{0.8cm}

\textit{Comprender las herramientas del derecho es el primer paso para crear valor real para quienes confían en un profesional.}

\vspace{1.2cm}

\rule{0.70\textwidth}{0.5pt}

\vspace{0.5cm}

{\large Apuntes de estudio}

\vspace{0.5cm}

\rule{0.70\textwidth}{0.5pt}

\vfill

{\large \autor}

\vspace{0.3cm}

{\large \anio}

\end{center}

\vfill

\begin{flushleft}
\textit{Este es un modesto aporte a los estudiantes de ``Derechos Reales'' no pretende sustituir el material oficial de la materia.}
\end{flushleft}

\end{titlepage}

% ── ÍNDICE ──────────────────────────────────────────────────
\tableofcontents
\thispagestyle{empty}
\newpage

% ============================================================
%  BLOQUE 1 · CASO INTRODUCTORIO Y MÉTODO DE ENCUADRE
% ============================================================
\chapter{Caso introductorio y método de encuadre}

\section{El inmueble desocupado de la familia}

\begin{example}{El inmueble desocupado de la familia}
Una familia es titular de varios bienes, entre ellos un inmueble que se encuentra totalmente desocupado. Esa misma familia necesita liquidez para reinvertirla en su empresa familiar. El interrogante planteado es el siguiente: \textbf{¿qué herramienta jurídica resulta recomendable antes de recurrir a la venta del bien?}
\end{example}

La primera reacción intuitiva ante un inmueble ocioso y una necesidad de liquidez suele ser la venta. Sin embargo, este enfoque no siempre resulta el más adecuado, especialmente cuando el bien tiene valor de legado o relevancia patrimonial para la familia. La distinción entre quien solo ofrece la compraventa como salida y quien acompaña a una familia a conservar su patrimonio explorando otras alternativas constituye una diferencia fundamental en el ejercicio profesional del derecho.

\begin{important}
El primer aprendizaje metodológico de esta unidad es la importancia de \textbf{encuadrar correctamente el problema} que plantea el cliente antes de ofrecer una solución. Si no se identifican con precisión las necesidades, las restricciones y los objetivos del cliente, es fácil quedarse en detalles secundarios y perder de vista lo que específicamente se requiere. Cuando alguien trae un problema, la primera tarea es cómo se lo encuadra.
\end{important}

Las herramientas jurídicas disponibles para un inmueble ocioso no se agotan en la venta. Entre las alternativas estudiadas en la materia se encuentran el usufructo, la locación y ---como se desarrolla en el bloque siguiente--- el derecho real de superficie, que permite generar ingresos sin desprenderse del dominio.

% ============================================================
%  BLOQUE 2 · FUNDAMENTOS DEL DERECHO REAL DE SUPERFICIE
% ============================================================
\chapter{Fundamentos del derecho real de superficie}

\section{Origen histórico del derecho de superficie}

El derecho real de superficie fue incorporado al ordenamiento jurídico argentino con la sanción del Código Civil y Comercial de la Nación (CCyC). Sin embargo, la figura ya contaba con un antecedente normativo previo: una \textbf{ley de forestación}, sancionada en un contexto en el que aún no se hablaba del derecho de superficie como categoría general, pero que respondía a la necesidad práctica de permitir forestar, sembrar y cosechar en terrenos ajenos sin necesidad de transferir la propiedad de esos terrenos.

\begin{example}{Por qué la forestación necesitaba esta herramienta}
La vida de un árbol no se mide en días sino en años: en muchos casos se necesitan entre 15 y 25 años para que una forestación dé sus frutos. Esa inversión no puede pensarse a corto plazo. Esta necesidad, sumada a las recomendaciones de organismos internacionales sobre el uso sostenible del medio ambiente, explica el origen de la figura: se requería una herramienta legal que le otorgara al forestador la seguridad de poder usar, gozar y disponer de lo plantado durante todos esos años, sin ser dueño del terreno.
\end{example}

\section{Concepto legal: artículo 2114 del CCyC}

Para comprender el derecho de superficie es necesario partir de la lógica clásica del \textbf{dominio}: cuando una persona es dueña de un terreno, es dueña del suelo, del subsuelo y del espacio aéreo por encima de él (el vuelo). El derecho de superficie rompe con ese paradigma, al permitir separar temporalmente esas tres capas y otorgarlas a distintos titulares.

\begin{formula}{Artículo 2114 CCyC --- Concepto}
\textit{``El derecho de superficie es un derecho real temporario, que se constituye sobre un inmueble ajeno, que otorga a su titular la facultad de uso, goce y disposición material y jurídica del derecho de plantar, forestar o construir, o sobre lo plantado, forestado o construido en el terreno, el vuelo o el subsuelo, según las modalidades de su ejercicio y plazo de duración establecidos en el título suficiente para su constitución, dentro de lo previsto en el presente Título y las leyes especiales.''}
\end{formula}

\begin{definition}{Derecho real de superficie}
El derecho real de superficie es un \textbf{derecho real temporario} que se constituye sobre un \textbf{inmueble ajeno} y otorga a su titular (el \textbf{superficiario}) la facultad de uso, goce y disposición material y jurídica para plantar, forestar o construir, o sobre lo ya plantado, forestado o construido, en el suelo, el vuelo o el subsuelo del inmueble, según las modalidades y el plazo pactados en el título constitutivo.
\end{definition}

Los dos elementos centrales de la definición son:

\begin{itemize}
    \item \textbf{Temporalidad:} a diferencia del dominio, que es perpetuo, el derecho de superficie tiene un plazo determinado.
    \item \textbf{Ajeneidad:} se constituye sobre un inmueble ajeno, mientras que el dominio recae sobre un inmueble propio.
\end{itemize}

\subsection{Diferencia con el dominio pleno}

En el dominio pleno, el propietario de un terreno es simultáneamente dueño del suelo, el subsuelo y el vuelo. Con el derecho de superficie, el propietario puede conservar la titularidad del suelo y otorgar a otra persona la titularidad del subsuelo o del vuelo en carácter de \textbf{dominio superficiario}: desmembrado y temporal.

\subsection{Diferencia con el usufructo}

En el dominio desmembrado por usufructo, el titular del dominio conserva ese derecho durante la vigencia del usufructo. En el derecho de superficie, el esquema es diferente: el superficiario tiene facultades de \textbf{uso, goce y disposición} sobre lo que construye, algo que el usufructuario no posee respecto de la cosa en sí misma.

\section{Objeto del derecho: plantar, forestar y construir}

El objeto del derecho de superficie comprende:

\begin{itemize}
    \item \textbf{Plantar, forestar o construir} sobre el suelo, el vuelo o el subsuelo ajenos.
    \item \textbf{Adquirir un derecho ya existente} sobre algo previamente plantado, forestado o construido: es decir, no es obligatorio partir de cero; puede constituirse un derecho de superficie sobre una construcción o plantación preexistente.
\end{itemize}

El derecho puede recaer sobre cualquiera de las tres capas del inmueble: el suelo, el subsuelo o el vuelo, según las modalidades pactadas en el título constitutivo.

\begin{example}{El estacionamiento bajo un edificio}
Un titular de terreno se reserva el subsuelo porque allí funciona un estacionamiento. El titular de ese garaje es dueño de él en carácter de \textbf{dominio superficiario}: desmembrado y temporal. Vencido el plazo, se producen las consecuencias reguladas en el artículo 2125 CCyC (véase la sección~\ref{sec:extincion_efectos}).
\end{example}

\section{El dominio superficiario temporal: pluralidad de superficiarios}

Un mismo inmueble puede estar afectado a \textbf{distintos derechos de superficie simultáneos} sobre distintas partes de él: el suelo, el subsuelo y el vuelo pueden entregarse a distintos superficiarios, cada uno con su propia habilitación y su propio plazo.

Esta posibilidad de superposición de derechos de superficie sobre un mismo inmueble es la que potencia el aprovechamiento económico del bien, como se ilustra en los casos prácticos del siguiente bloque.

% ============================================================
%  BLOQUE 3 · CASOS PRÁCTICOS DE APLICACIÓN
% ============================================================
\chapter{Casos prácticos de aplicación}

\section{Múltiples derechos de superficie sobre un mismo terreno}

\begin{example}{El terreno frente al predio ferial}
La titular de un terreno ubicado frente a un predio ferial, adquirido en subasta pública, no tiene intención de invertir en construir porque toda su inversión se agotó en la compra. Ante esa situación, aparecen distintos interesados:

\begin{enumerate}
    \item Un interesado propone construir un garaje o estacionamiento en el subsuelo. Se le otorga un \textbf{derecho de superficie por quince años sobre el subsuelo}.
    \item Un segundo interesado quiere desarrollar un centro cultural en la planta. Se le otorga un \textbf{derecho de superficie por veinte años}.
    \item Un tercer interesado propone construir, por encima del centro cultural (desde el segundo piso hacia arriba), un centro de convenciones. Se le cede por \textbf{veinte años la porción del derecho de superficie correspondiente al vuelo}.
\end{enumerate}

Cada superficiario invierte en su propia obra y la explota durante el plazo convenido. La titular del terreno cobra retribución de los tres. El dominio nunca se transfiere: al vencimiento de cada plazo, los bienes construidos revierten al patrimonio de la titular.
\end{example}

Este mecanismo permite \textbf{aprovechar un mismo terreno en varios niveles y con varios destinos simultáneamente} sin que la propietaria pierda la titularidad del dominio ni deba vender el bien.

\section{El convenio con el gobierno de la ciudad: el cargo como herramienta de protección}

En el contexto de un convenio entre el gobierno de la ciudad y el colegio profesional para la escrituración de un barrio construido para personas provenientes de un asentamiento informal, algunos predios estaban destinados a asignarse a asociaciones civiles para su funcionamiento.

La cuestión jurídica planteada fue: ¿conviene transferir el dominio pleno a la asociación civil, o existe una herramienta más adecuada para proteger el destino social del bien?

El riesgo de la donación en dominio pleno es evidente: si la asociación civil deja de realizar los actos para los cuales le fue asignado el inmueble, o si eventualmente lo vende, el objetivo perseguido por el gobierno al asignarlo se pierde definitivamente.

\begin{example}{La escritura con cargo}
La solución propuesta fue constituir un \textbf{derecho de superficie con cargo}: el derecho se extingue si no se cumple el objetivo para el cual fue otorgado, es decir, si la asociación civil deja de desarrollar la actividad comprometida. En ese caso, el inmueble puede reasignarse a otra asociación civil con interés en desarrollar esa actividad, sin que el predio quede desocupado ni se use en beneficio de intereses ajenos al objetivo original.
\end{example}

Este mecanismo garantiza que el derecho proteja la \textbf{finalidad social del bien público} más allá del acto formal de la escrituración.

\section{El legado familiar y la alternativa a la donación}

Cuando una familia posee un terreno con \textbf{valor de legado} ---por ejemplo, el terreno donde nació y vivió parte de la familia durante generaciones--- pero ya no puede mantenerlo económicamente, la alternativa a la donación o venta en dominio pleno es constituir un \textbf{derecho de superficie con cargo}.

\begin{example}{La donación del abuelo}
Un titular había comprado un terreno con la intención de destinarlo, junto con terrenos linderos, al desarrollo de un club social. Todos los propietarios donaron sus terrenos en dominio pleno. Sin embargo, el club nunca se desarrolló como estaba previsto, y ese terreno terminó siendo vendido para otra obra de bien público, completamente ajena al proyecto original que había motivado la donación.

Si en lugar de donar el dominio pleno se hubiera constituido un derecho de superficie con cargo, el terreno habría podido volver a los donantes (o a sus herederos) al no cumplirse el objetivo pactado, preservando así el legado familiar.
\end{example}

La conclusión práctica es que, en muchas situaciones, \textbf{en lugar de transferir el dominio pleno conviene constituir un derecho de superficie con cargo}, de modo que el bien vuelva a la familia si no se cumple la finalidad pactada.

\section{La revalorización de zona: el caso del hipermercado}

Cuando una gran cadena comercial se instala en una zona, genera a su alrededor un polo de actividad: galerías, locales y otros comercios. Pero esa cadena puede no querer comprar el terreno de manera definitiva, porque no tiene certeza sobre el rendimiento de esa ubicación y no desea hacer una inversión tan grande en ese momento.

\begin{example}{Por qué el hipermercado prefiere un derecho de superficie}
Mediante un derecho de superficie por veinte años, la cadena construye el hipermercado, las galerías y el estacionamiento, y explota el negocio durante ese plazo. Al vencimiento, puede renovar el derecho, asociarse con el titular del terreno o pagarle un alquiler por continuar.

Para el titular del terreno, este mecanismo resulta doblemente beneficioso: cobra la retribución pactada con el superficiario y, además, si posee inmuebles en la zona, su valor aumenta considerablemente por el solo hecho de tener ese polo comercial cerca.
\end{example}

La ventaja adicional frente al usufructo es que el titular del terreno \textbf{no pierde el dominio}: lo cede temporalmente y, al vencimiento, recupera el bien junto con las mejoras realizadas por el superficiario.

\section{Superficie sin ser dueño del suelo: el agronegocio y los bonos verdes}

En el ámbito agropecuario, empresas agroindustriales pueden desarrollar su actividad sin ser propietarias de los terrenos que explotan, aprovechando el derecho real de superficie. Esto les permite trasladarse de una zona climática a otra y cosechar a gran escala sin disponer del capital necesario para adquirir los terrenos.

Las alternativas tradicionales a este esquema son los contratos de arrendamiento, que presentan desventajas para ambas partes: el arrendatario no tiene la seguridad de un derecho real inscribible, y el arrendador puede ver limitada su capacidad de renegociación.

El derecho de superficie también se vincula con el mercado de \textbf{bonos verdes}: las empresas que emiten dióxido de carbono pueden compensar ese impacto promoviendo bosques forestales a través de la compra de bonos verdes. Esta práctica, muy extendida en Europa, podría aplicarse en Argentina mediante derechos de superficie orientados a la forestación, con beneficios directos para los parques nacionales y para el desarrollo sostenible.

% ============================================================
%  BLOQUE 4 · EXTINCIÓN, REGISTRO Y LIQUIDACIÓN
% ============================================================
\chapter{Extinción, registro y liquidación}

\section{Extinción del derecho de superficie}
\label{sec:extincion}

La extinción del derecho de superficie está regulada en el artículo 2124 del CCyC.

\begin{formula}{Artículo 2124 CCyC --- Extinción}
\textit{``El derecho de construir, plantar o forestar se extingue por renuncia expresa, vencimiento del plazo, cumplimiento de una condición resolutoria, por consolidación y por el no uso durante diez años, para el derecho a construir, y de cinco años, para el derecho a plantar o forestar.''}
\end{formula}

Las causales de extinción son, en consecuencia:

\begin{itemize}
    \item \textbf{Renuncia expresa} del superficiario.
    \item \textbf{Vencimiento del plazo} convenido en el título constitutivo.
    \item \textbf{Cumplimiento de una condición resolutoria} pactada.
    \item \textbf{Consolidación} (reunión en una misma persona de la titularidad del suelo y del derecho de superficie).
    \item \textbf{No uso:} diez años para el derecho a construir; cinco años para el derecho a plantar o forestar.
\end{itemize}

Las partes pueden pactar en la escritura constitutiva \textbf{causales adicionales de extinción}, por ejemplo, limitando el tipo de actividad que el superficiario puede desarrollar.

\begin{important}
Si bien no se puede prohibir directamente la transmisión del derecho de superficie, sí es posible \textbf{limitarla}: por ejemplo, pactando que el derecho sea cedido a terceros solamente en tanto y en cuanto estos no desarrollen la misma actividad que desarrolla el titular del suelo (para evitar que, mediante la cesión, termine instalándose un competidor directo).
\end{important}

\section{Efectos de la extinción}
\label{sec:extincion_efectos}

Los efectos de la extinción están regulados en el artículo 2125 del CCyC.

\begin{formula}{Artículo 2125 CCyC --- Efectos de la extinción}
\textit{``Al momento de la extinción del derecho de superficie por el cumplimiento del plazo convencional o legal, el propietario del suelo hace suyo lo construido, plantado o forestado, libre de los derechos reales o personales impuestos por el superficiario. Si el derecho de superficie se extingue antes del cumplimiento del plazo legal o convencional, los derechos reales constituidos sobre la superficie o sobre el suelo continúan gravando separadamente las dos partes, como si no hubiese habido extinción, hasta el transcurso del plazo del derecho de superficie. Subsisten también los derechos personales durante el tiempo establecido.''}
\end{formula}

En síntesis:

\begin{itemize}
    \item Si la extinción se produce por \textbf{vencimiento del plazo}: el propietario del suelo hace suyo lo construido, plantado o forestado, \textbf{libre de las cargas} impuestas por el superficiario.
    \item Si la extinción se produce \textbf{antes del plazo pactado}: los derechos reales y personales constituidos sobre la superficie continúan vigentes hasta el vencimiento del plazo original, como si la extinción no se hubiera producido.
\end{itemize}

\section{Modalidades de pago y liquidación final}

La retribución al titular del suelo por constituir el derecho de superficie puede pactarse de distintas maneras:

\begin{enumerate}
    \item \textbf{Pago total al contado} en el momento de suscribir el derecho de superficie.
    \item \textbf{Pago en cuotas o plazos}, a medida que se van percibiendo los ingresos de la explotación (por ejemplo, de alquileres temporarios).
    \item \textbf{Cesión de un porcentaje de las unidades construidas} en favor del titular del suelo como retribución (por ejemplo, el diez por ciento de los departamentos construidos).
\end{enumerate}

Al momento de la extinción, el superficiario \textbf{no está obligado a demoler} lo construido y devolver el terreno vacío: lo entrega junto con la construcción realizada. También es posible pactar de antemano cómo se reparten las unidades construidas al finalizar el plazo.

\begin{example}{Liquidación al vencimiento}
Si se construyeron veinte departamentos, puede pactarse que al vencer el plazo el superficiario conserve ocho y el titular del suelo reciba doce, extinguiéndose el derecho de superficie y transmitiéndose la propiedad de las unidades según lo convenido. El superficiario hace suyo el dominio de la parte de lo construido que le corresponde según el acuerdo.
\end{example}

\section{Aspectos registrales}

La inscripción del derecho de superficie se realiza mediante \textbf{escritura pública}, que se inscribe en el Registro de la Propiedad dentro de la \textbf{misma matrícula del inmueble}, en la parte correspondiente al dominio, del mismo modo en que se inscribe un usufructo.

\begin{itemize}
    \item La \textbf{titularidad del dominio} permanece a nombre del propietario del suelo.
    \item El \textbf{derecho de superficie} se anota en la matrícula como gravamen, con la misma minuta con la que se anota un dominio.
\end{itemize}

La \textbf{cancelación registral} al finalizar el plazo es análoga a la cancelación de una hipoteca: del mismo modo en que una hipoteca se cancela al cumplirse su objeto, el derecho de superficie se cancela y el derecho real registrado se extingue. En todos los casos, la cancelación registral debe contemplar lo pactado en la escritura, incluyendo eventuales pactos complementarios dentro del mismo instrumento.

\begin{example}{La cochera bajo la iglesia}
Una iglesia situada sobre una loma no deseaba vender el terreno bajo ningún concepto, dado que sobre él se encuentra el templo. Para aprovechar el subsuelo sin desprenderse del dominio, se constituyeron distintos derechos de superficie que permitieron desarrollar una cochera debajo del edificio religioso, sin afectar la titularidad de la iglesia sobre el inmueble.
\end{example}

% ============================================================
%  BLOQUE 5 · DIMENSIÓN ECONÓMICA Y COMPARACIÓN FINAL
% ============================================================
\chapter{Dimensión económica y comparación final}

\section{Superficie, locación y usufructo: ventajas comparativas}

\subsection{Comparación con la locación (contrato de alquiler)}

El principal argumento a favor del derecho de superficie frente a la locación es de naturaleza jurídica: \textbf{el derecho de superficie es un derecho real}, oponible erga omnes, mientras que la locación es un derecho personal (contractual).

\begin{example}{El depósito de camionetas}
Una empresa propone instalar un depósito de camionetas en un terreno familiar heredado con valor sentimental. La titular no quiere vender, pero la empresa necesita invertir sumas importantes en seguridad y tecnología para el depósito. Si se celebra un simple contrato de locación, esa inversión queda a merced de la duración del contrato y de la voluntad de las partes al vencimiento, lo que genera incertidumbre sobre la recuperación de lo invertido.

El derecho de superficie, en cambio, le otorga a la empresa un \textbf{derecho real inscribible, con plazo cierto}, sobre el cual puede invertir con seguridad jurídica. Al finalizar el plazo, la titular podría incluso vender ese derecho de superficie a un tercero.
\end{example}

\vspace{0.5em}

\begin{table}[H]
\centering
\caption{Comparación entre el derecho de superficie y la locación}
\begin{tabularx}{\textwidth}{
    >{\raggedright\arraybackslash}p{0.30\textwidth}
    >{\raggedright\arraybackslash}X
    >{\raggedright\arraybackslash}X
}
\toprule
\textbf{Criterio} & \textbf{Derecho de superficie} & \textbf{Locación (alquiler)} \\
\midrule
Naturaleza jurídica & Derecho real, oponible erga omnes & Derecho personal (contrato) \\
\addlinespace
Inscripción registral & Sí, en la matrícula del inmueble & No se inscribe como derecho real \\
\addlinespace
Facultad de disposición & El superficiario puede disponer del derecho & El locatario no puede disponer de la cosa \\
\addlinespace
Posibilidad de modificar la cosa & Sí, es la esencia del derecho (construir, plantar, forestar) & En principio no, salvo autorización expresa \\
\addlinespace
Embargabilidad & El derecho es embargable & El contrato en sí no es un bien embargable del mismo modo \\
\addlinespace
Seguridad para quien invierte & Alta: derecho real con plazo cierto & Baja: depende de la duración y renovación del contrato \\
\addlinespace
Transmisión por muerte del titular & Se transmite a los herederos & Depende de lo pactado y de la ley aplicable al contrato \\
\bottomrule
\end{tabularx}
\end{table}

\subsection{Comparación con el usufructo}

La diferencia más relevante entre el derecho de superficie y el usufructo es la siguiente:

\begin{itemize}
    \item \textbf{Usufructo:} se extingue con la muerte del usufructuario. No se transmite a los herederos.
    \item \textbf{Derecho de superficie:} no se extingue por el fallecimiento del superficiario. Se transmite a sus herederos.
\end{itemize}

Cuando el superficiario es una \textbf{persona jurídica}, existe además una ventaja operativa adicional: en lugar de transmitir el derecho de superficie mediante una nueva escritura, puede cederse la participación societaria (acciones o cuotas de esa persona jurídica), simplificando considerablemente la operación.

\begin{important}
Esta facilidad de transferencia mediante cesión de acciones conlleva un riesgo: si se contrató con una empresa por su reputación específica, y luego esa empresa cede sus acciones a otra persona sin la misma solvencia o reputación, el titular del terreno puede terminar vinculado, de hecho, con una contraparte distinta de la que originalmente eligió. Esta circunstancia debe considerarse al redactar la escritura constitutiva.
\end{important}

\section{El efecto multiplicador económico y social}

La actividad de la construcción no se agota en la obra en sí misma: genera una \textbf{cadena de actividades conexas} que derrama beneficios hacia otros sectores. La obra da trabajo al arquitecto, a los obreros y a los distintos gremios; esos trabajadores, a su vez, dinamizan el comercio local. Fomentar el derecho real de superficie genera, en consecuencia, un efecto multiplicador que excede ampliamente el negocio inmobiliario puntual.

\begin{example}{El local con dos pisos para coworking}
Un comerciante es dueño de un local y recibe la propuesta de que alguien construya dos pisos por encima para destinarlos a un espacio de coworking. El comerciante teme comprometer el destino futuro de su esquina, que podría valorizarse en los próximos años.

La solución es el derecho de superficie: el comerciante sigue explotando su local en planta baja, cede por un plazo determinado el derecho de construir y explotar el coworking en los pisos superiores, y al vencimiento del plazo recupera la disponibilidad plena del inmueble, incluidas las mejoras realizadas por el superficiario. Quien invierte en la construcción tiene la seguridad jurídica necesaria para hacerlo ---puede incluso vender su derecho durante la vigencia del plazo--- y el propietario original no compromete el destino futuro de su propiedad.
\end{example}

Esta misma lógica puede aplicarse al \textbf{problema habitacional}: en lugar de transferir el dominio pleno de una vivienda, pueden otorgarse derechos de uso temporal por plazos determinados ---por ejemplo, diez años--- a jóvenes que están comenzando su vida económica. Transcurrido ese plazo, esas personas estarían en mejores condiciones de acceder a un crédito y a una vivienda propia, liberando la unidad para que otra persona joven pueda utilizarla en las mismas condiciones.

\begin{example}{El modelo de vivienda comunitaria en los kibutz de Israel}
En los kibutz de Israel existen unidades habitacionales de dos, tres y cuatro ambientes que se asignan según la etapa de vida de cada familia: una pareja recién casada comienza en una unidad pequeña; cuando llegan los hijos, se muda a una mayor; cuando esos hijos crecen y se van, la pareja vuelve a una unidad más pequeña, liberando la mayor para otra familia con hijos. La propiedad no es individual: la vivienda se comparte entre la comunidad. De ese modo, el espacio habitacional se aprovecha de manera constante, sin quedar ocioso ni concentrado en pocas manos durante toda la vida de una sola familia.
\end{example}

La resistencia cultural a utilizar estas herramientas no se debe a la falta de herramientas legales disponibles, sino a la \textbf{falta de costumbre y de creatividad} tanto de los clientes como de los propios profesionales que los asesoran. El derecho es una herramienta de creación cuando se lo ejerce con imaginación.

% ============================================================
%  CUADRO CORNELL
% ============================================================
\chapter{Cuadro Cornell: repaso integral}

El siguiente cuadro sintetiza los conceptos centrales de la materia en formato Cornell. La columna izquierda contiene las preguntas clave; la columna derecha, el desarrollo correspondiente. La última fila presenta la síntesis integradora.

\vspace{0.5em}

\begin{table}[H]
\centering
\begin{tabularx}{\textwidth}{
    >{\raggedright\arraybackslash}p{0.27\textwidth}
    >{\raggedright\arraybackslash}X
}
\toprule
\textbf{Preguntas / palabras clave} & \textbf{Notas / respuestas} \\
\midrule

\textbf{¿Qué es el derecho real de superficie?} &
Derecho real temporario (art.~2114 CCyC) constituido sobre un inmueble ajeno, que otorga al superficiario la facultad de uso, goce y disposición material y jurídica para plantar, forestar o construir, o sobre lo ya plantado, forestado o construido, en el suelo, el vuelo o el subsuelo del inmueble, según las modalidades y el plazo pactados. \\
\addlinespace

\textbf{¿Cuál fue su antecedente histórico?} &
Una ley de forestación, anterior al CCyC, pensada para permitir plantar, forestar o cosechar en terrenos ajenos sin transferir su propiedad, dado que los ciclos forestales requieren inversiones de muy largo plazo (entre 15 y 25 años). \\
\addlinespace

\textbf{¿Qué paradigma del dominio rompe?} &
El paradigma de que el dueño del terreno es simultáneamente dueño del suelo, el subsuelo y el vuelo. Con la superficie, esas tres capas pueden separarse y otorgarse a distintos titulares, cada uno dueño superficiario de su porción, por tiempo determinado. \\
\addlinespace

\textbf{¿Puede haber más de un derecho de superficie sobre el mismo inmueble?} &
Sí. Un mismo dominio puede estar afectado a distintos derechos superficiales simultáneos sobre el suelo, el vuelo y el subsuelo (ejemplo: estacionamiento en el subsuelo, centro cultural en planta y centro de convenciones en los pisos superiores). \\
\addlinespace

\textbf{¿Cómo protege este derecho el destino social o familiar de un bien?} &
Mediante la constitución de un derecho de superficie \textbf{con cargo}: si el superficiario no cumple la finalidad pactada, el derecho se extingue y el bien puede reasignarse. A diferencia de una donación en dominio pleno, esta figura ofrece reversión ante el incumplimiento. \\
\addlinespace

\textbf{¿Cuándo se extingue el derecho de superficie? (art.~2124)} &
Por renuncia expresa, vencimiento del plazo, cumplimiento de una condición resolutoria, consolidación, y por el no uso durante diez años (derecho a construir) o cinco años (derecho a plantar o forestar). Las partes pueden pactar causales adicionales. \\
\addlinespace

\textbf{¿Qué ocurre con lo construido al extinguirse el derecho? (art.~2125)} &
Si se extingue por vencimiento del plazo: el propietario del suelo hace suyo lo construido, libre de cargas. Si se extingue antes del plazo, los derechos reales y personales constituidos sobre la superficie continúan vigentes hasta el vencimiento del plazo original. \\
\addlinespace

\textbf{¿Cómo se retribuye al titular del suelo?} &
De tres maneras posibles según lo pactado: pago total al contado al constituir el derecho; pago en cuotas a medida que se generan ingresos; o cesión de un porcentaje de las unidades construidas (p.~ej., el 10\% de los departamentos). \\
\addlinespace

\textbf{¿Cómo se inscribe registralmente?} &
Mediante escritura pública, inscripta en la misma matrícula del inmueble, en la parte correspondiente al dominio, del mismo modo que se inscribe un usufructo. La titularidad del dominio permanece a nombre del propietario del suelo; el derecho de superficie se anota como gravamen. \\
\addlinespace

\textbf{¿Qué ventajas tiene frente a la locación?} &
Es un derecho real: se inscribe, es embargable, permite modificar la cosa (construir, plantar, forestar) y da seguridad jurídica a largo plazo a quien invierte. Un contrato de locación no ofrece las mismas garantías para inversiones de largo plazo. \\
\addlinespace

\textbf{¿Qué ventajas tiene frente al usufructo?} &
No se extingue con la muerte del superficiario: se transmite a sus herederos. El usufructo, en cambio, se extingue con la muerte del usufructuario. Cuando el superficiario es persona jurídica, el derecho puede transferirse de modo simplificado cediendo las acciones o cuotas de esa sociedad. \\
\addlinespace

\textbf{¿Qué impacto económico y social tiene?} &
Genera un efecto multiplicador: activa la construcción, da trabajo a arquitectos, obreros y gremios, y dinamiza el comercio local. Permite además soluciones creativas al problema habitacional y al uso eficiente de terrenos del Estado o de instituciones, sin necesidad de vender o donar el dominio pleno. \\

\midrule

\multicolumn{2}{p{\textwidth}}{%
\textbf{Síntesis:} El derecho real de superficie ---regulado desde el art.~2114 CCyC y con antecedente en la ley de forestación--- permite separar temporalmente el suelo, el vuelo y el subsuelo de un inmueble para otorgar a otra persona, física o jurídica, la facultad de plantar, forestar o construir sobre ellos por un plazo determinado, sin transferir el dominio. A través de múltiples casos prácticos (el terreno frente al predio ferial con tres superficiarios distintos, el convenio con el gobierno para asociaciones civiles, el legado familiar, el hipermercado que revaloriza una zona, el agronegocio sin tierra propia, el depósito de camionetas y el local con coworking) se evidencia que este derecho ofrece seguridad jurídica al inversor (a diferencia de la locación), se transmite por muerte a los herederos (a diferencia del usufructo), admite pactos flexibles de pago y reversión mediante cargos, y genera un efecto multiplicador económico y social que excede el negocio inmobiliario puntual.%
} \\

\bottomrule
\end{tabularx}
\end{table}

\end{document}
